% 康奈尔听课笔记模板（配色取自王道课件）
\documentclass[10.5pt,a4paper,oneside]{article}
\usepackage[a4paper,top=2.3cm,bottom=2.2cm,left=5.7cm,right=1.5cm,
  marginparwidth=4.25cm,marginparsep=0.55cm,headheight=14pt,footskip=1.1cm]{geometry}
\usepackage{xeCJK}
\usepackage{xeCJKfntef}
\xeCJKDeclareCharClass{FullLeft}{"201C,"2018}
\xeCJKDeclareCharClass{FullRight}{"201D,"2019}
\xeCJKDeclareCharClass{CJK}{"2460->"24FF,"2605,"2606}
\setmainfont{Noto Sans CJK SC}[BoldFont={Noto Sans CJK SC Bold},ItalicFont={Noto Sans CJK SC},ItalicFeatures={FakeSlant=0.18}]
\setsansfont{Noto Sans CJK SC}[BoldFont={Noto Sans CJK SC Bold}]
\setmonofont{DejaVu Sans Mono}[Scale=0.9]
\setCJKmainfont{Noto Sans CJK SC}[BoldFont={Noto Sans CJK SC Bold},ItalicFont={Noto Sans CJK SC},ItalicFeatures={FakeSlant=0.18}]
\setCJKsansfont{Noto Sans CJK SC}[BoldFont={Noto Sans CJK SC Bold}]
\setCJKmonofont{Noto Sans Mono CJK SC}[Scale=0.9]
\linespread{1.32}
\usepackage{amsmath,amssymb}
\usepackage[dvipsnames,table]{xcolor}
\usepackage{graphicx}
\usepackage{adjustbox}
\usepackage{tikz}
\usetikzlibrary{calc}
\usepackage[most]{tcolorbox}
\usepackage{eso-pic}
\usepackage{fancyhdr}
\usepackage{titlesec}
\usepackage{needspace}
\usepackage{enumitem}
\usepackage{xltabular,booktabs,array,colortbl}
\usepackage{marginnote}
\reversemarginpar
\renewcommand*{\raggedleftmarginnote}{}
\renewcommand*{\raggedrightmarginnote}{}
\usepackage{caption}
\usepackage{titletoc}
\usepackage{fancyvrb}
\usepackage{framed}
\usepackage[section]{placeins}
\usepackage{hyperref}
\hypersetup{colorlinks=true,linkcolor=wdblue,urlcolor=wdblue,pdftitle={2.2.1
顺序表的定义},pdfauthor={听课笔记}}

% ---------- 课件配色 ----------
\definecolor{wdblue}{HTML}{4472C4}
\definecolor{wdlightblue}{HTML}{5B9BD5}
\definecolor{wdtitle}{HTML}{4A86C8}
\definecolor{wddark}{HTML}{1F3864}
\definecolor{wdorange}{HTML}{ED7D31}
\definecolor{wdgray}{HTML}{A5A5A5}
\definecolor{wdgold}{HTML}{FFC000}
\definecolor{wdgreen}{HTML}{70AD47}
\definecolor{wddeepgreen}{HTML}{3F7F2A}
\definecolor{wdred}{HTML}{C00000}
\definecolor{wdpurple}{HTML}{7030A0}
\definecolor{cuebg}{HTML}{F2F6FC}
\definecolor{codebg}{HTML}{F7F9FC}
\color[HTML]{262626}

% ---------- 行内样式 ----------
\newcommand{\hlred}[1]{\textcolor{wdred}{\bfseries #1}}
\newcommand{\hlblue}[1]{\textcolor{wdblue}{\bfseries #1}}
\newcommand{\hlgreen}[1]{\textcolor{wddeepgreen}{\bfseries #1}}
\newcommand{\hlorange}[1]{\textcolor{wdorange}{\bfseries #1}}
\newcommand{\hlpurple}[1]{\textcolor{wdpurple}{\bfseries #1}}
\newcommand{\uln}[1]{\CJKunderline[format=\color{wdred},thickness=0.9pt,depth=0.28em]{#1}}
\newcommand{\wavy}[1]{\CJKunderwave[format=\color{wdorange}]{#1}}
\newcommand{\hlmark}[1]{\CJKunderline[format=\color{wdgold!55},thickness=0.55em,depth=-0.05em]{#1}}
\newcommand{\tagbox}[2]{\tcbox[on line,colback=#1,colframe=#1,coltext=white,boxsep=0.5pt,left=3pt,right=3pt,top=0.5pt,bottom=0.5pt,arc=2pt,fontupper=\small\bfseries]{#2}}

% ---------- 康奈尔版式背景：左侧线索栏 ----------
\newlength{\cuestripw}\setlength{\cuestripw}{5.15cm}
\AddToShipoutPictureBG{%
  \begin{tikzpicture}[remember picture,overlay]
    \fill[cuebg] (current page.north west) rectangle ($(current page.south west)+(\cuestripw,0)$);
    \draw[wdblue,line width=1.2pt] ($(current page.north west)+(\cuestripw,-1.55cm)$) -- ($(current page.south west)+(\cuestripw,1.55cm)$);
    \node[anchor=west,text=wdblue,font=\small\bfseries] at ($(current page.north west)+(0.75cm,-1.25cm)$) {线索 · 问题};
    \node[anchor=west,text=wdblue,font=\small\bfseries] at ($(current page.north west)+(\cuestripw+0.55cm,-1.25cm)$) {笔记 · 老师主线 + 补充讲解};
  \end{tikzpicture}}

% ---------- 页眉页脚 ----------
\pagestyle{fancy}\fancyhf{}
\renewcommand{\headrulewidth}{0pt}
\fancyhead[R]{\small\color{wdgray}第2章 线性表}
\fancyfoot[C]{\small\color{wdgray}\thepage}
\fancyfoot[R]{\small\color{wdorange}2.2.1 顺序表的定义 · 听课笔记}
\fancyfoot[L]{\small\hyperlink{toc}{\textcolor{wdblue}{\ensuremath{\uparrow}\,回到目录}}}
\fancypagestyle{plain}{\fancyhf{}\fancyfoot[C]{\small\color{wdgray}\thepage}\fancyfoot[R]{\small\color{wdorange}2.2.1
顺序表的定义 · 听课笔记}}

% ---------- 标题：仿课件的蓝色标题 + 六色条 ----------
\newcommand{\sixbar}{\par\nointerlineskip\vspace{3pt}\centering
  \begin{tikzpicture}\foreach \c [count=\i] in {wdlightblue,wdorange,wdgray,wdgold,wdblue,wdgreen}
     {\fill[\c] ({(\i-1)*1.05},0) rectangle ({(\i-1)*1.05+0.92},0.11);}\end{tikzpicture}\par}
\titleformat{\section}[block]{\centering\Large\bfseries\color{wdtitle}}{}{0pt}{}[\sixbar]
\titlespacing*{\section}{0pt}{16pt}{10pt}
\pretocmd{\section}{\Needspace{8\baselineskip}}{}{}
\titleformat{\subsection}{\Needspace{7\baselineskip}\large\bfseries\color{wddark}}{\textcolor{wdblue}{\rule[-0.15em]{4pt}{1.05em}}\hspace{6pt}}{0pt}{}
\titlespacing*{\subsection}{0pt}{12pt}{5pt}
\titleformat{\subsubsection}{\Needspace{5\baselineskip}\normalsize\bfseries\color{wdblue}}{\textcolor{wdorange}{\large\ensuremath{\blacktriangleright}}\hspace{4pt}}{0pt}{}
\titlespacing*{\subsubsection}{0pt}{8pt}{3pt}
\setcounter{secnumdepth}{0}
\setlength{\parindent}{0pt}\setlength{\parskip}{5pt}
\setlist{nosep,leftmargin=1.5em,itemsep=2pt}
\setlist[itemize,1]{label=\textcolor{wdblue}{\small\ensuremath{\bullet}}}
\setlist[itemize,2]{label=\textcolor{wdorange}{\scriptsize\ensuremath{\blacktriangleright}}}
\setlist[enumerate,1]{label=\textcolor{wdblue}{\bfseries\arabic*.}}
\providecommand{\tightlist}{\setlength{\itemsep}{2pt}\setlength{\parskip}{0pt}}

% ---------- 目录 ----------
\setcounter{tocdepth}{1}
\titlecontents{section}[0pt]{\addvspace{5pt}\bfseries\color{wddark}}{}{}{\titlerule*[0.7pc]{\color{wdgray!60}.}\contentspage}
\titlecontents{subsection}[1.6em]{\addvspace{1pt}\small\color{black!75}}{}{}{\titlerule*[0.7pc]{\color{wdgray!40}.}\contentspage}
\newcommand{\tocback}{\par\vspace{-3pt}{\hfill\footnotesize\hyperlink{toc}{\textcolor{wdgray}{\ensuremath{\uparrow}\,回到目录}}}\par\vspace{-4pt}}

% ---------- 线索栏 ----------
\newcommand{\cueblock}[2]{\Needspace{4\baselineskip}\leavevmode#1{#2}\par\vspace{-\baselineskip}\vspace{-\parskip}}
% 线索直接贴在段首行的左侧（不用 marginnote，避免分页时偶尔落进正文栏）
\newcommand{\cuebox}[1]{\leavevmode\llap{\smash{\begin{minipage}[t]{\marginparwidth}\raggedright\small\linespread{1.18}\selectfont #1\end{minipage}}\hspace{\marginparsep}}}
\newcommand{\cuenote}[1]{\cuebox{{\color{wdblue}\bfseries Q\,}\color{wddark}#1}}
\newcommand{\cuekey}[1]{\cuebox{\color{wdred}\bfseries #1}}

% ---------- 讲解框 ----------
\tcbset{wdbox/.style={enhanced,breakable,arc=2pt,boxrule=0pt,leftrule=3pt,left=7pt,right=7pt,top=5pt,bottom=5pt,
   before skip=7pt,after skip=7pt,fonttitle=\bfseries\small,attach title to upper={\quad},coltitle=#1,
   colframe=#1,colback=#1!6}}
\newtcolorbox{defbox}[1][]{enhanced,breakable,arc=3pt,boxrule=0.8pt,colframe=wdblue,colback=wdblue!4,
   left=8pt,right=8pt,top=6pt,bottom=6pt,fonttitle=\bfseries\small,coltitle=white,colbacktitle=wdblue,
   title={定义#1},before skip=8pt,after skip=8pt}
\newtcolorbox{bridgebox}{wdbox=wdblue,title={承上启下}}
\newtcolorbox{plainbox}{wdbox=wddeepgreen,colframe=wdgreen,colback=wdgreen!8,title={大白话}}
\newtcolorbox{deepbox}{wdbox=wdpurple,title={背后的思想}}
\newtcolorbox{exambox}{wdbox=wdred,title={考点}}
\newtcolorbox{pitfallbox}{wdbox=wdorange,title={易错提醒}}
\newtcolorbox{bookbox}{wdbox=wddark,colback=wdgray!10,colframe=wdgray,title={王道书补充}}
\newtcolorbox{teacherbox}{wdbox=wdorange,colback=wdgold!8,colframe=wdgold,coltitle=wdorange,title={老师原话}}
\newtcolorbox{fixbox}{wdbox=wdred,colback=wdred!4,title={勘误 / 说明}}
% 习题笔记用
\newtcolorbox{questionbox}[1][]{enhanced,breakable,arc=3pt,boxrule=0.8pt,colframe=wddark,colback=white,
   left=8pt,right=8pt,top=6pt,bottom=6pt,fonttitle=\bfseries\small,coltitle=white,colbacktitle=wddark,
   title={#1},before skip=10pt,after skip=6pt}
\newtcolorbox{videobox}{wdbox=wdorange,colback=wdgold!8,colframe=wdgold,coltitle=wdorange,title={视频讲解}}
\newtcolorbox{jiexibox}{wdbox=wddark,colback=wdgray!10,colframe=wdgray,title={书上解析}}
\newtcolorbox{addbox}{wdbox=wdpurple,title={我的补充}}
\newtcolorbox{recallbox}{wdbox=wdblue,title={回顾正课}}
\newtcolorbox{watchbox}{enhanced,breakable,arc=3pt,boxrule=1pt,colframe=wdorange,colback=wdgold!12,
   fonttitle=\bfseries\small,coltitle=white,colbacktitle=wdorange,title={建议去看视频},left=8pt,right=8pt,before skip=8pt,after skip=8pt}
\newtcolorbox{summarybox}{enhanced,arc=3pt,
   boxrule=1pt,colframe=wdblue,colback=white,colbacktitle=wdblue,coltitle=white,fonttitle=\bfseries,
   title={总结区 · 用自己的话复述这一节},left=10pt,right=10pt,top=8pt,bottom=8pt,before skip=14pt}

% ---------- 代码 ----------
\usepackage{color}
\usepackage{fancyvrb}
\newcommand{\VerbBar}{|}
\newcommand{\VERB}{\Verb[commandchars=\\\{\}]}
\DefineVerbatimEnvironment{Highlighting}{Verbatim}{commandchars=\\\{\}}
% Add ',fontsize=\small' for more characters per line
\newenvironment{Shaded}{}{}
\newcommand{\AlertTok}[1]{\textcolor[rgb]{1.00,0.00,0.00}{#1}}
\newcommand{\AnnotationTok}[1]{\textcolor[rgb]{0.38,0.63,0.69}{#1}}
\newcommand{\AttributeTok}[1]{\textcolor[rgb]{0.49,0.56,0.16}{#1}}
\newcommand{\BaseNTok}[1]{\textcolor[rgb]{0.77,0.35,0.07}{#1}}
\newcommand{\BuiltInTok}[1]{\textcolor[rgb]{0.00,0.50,0.00}{#1}}
\newcommand{\CharTok}[1]{\textcolor[rgb]{0.63,0.19,0.43}{#1}}
\newcommand{\CommentTok}[1]{\textcolor[rgb]{0.18,0.55,0.23}{#1}}
\newcommand{\CommentVarTok}[1]{\textcolor[rgb]{0.18,0.55,0.23}{#1}}
\newcommand{\ConstantTok}[1]{\textcolor[rgb]{0.77,0.35,0.07}{#1}}
\newcommand{\ControlFlowTok}[1]{\textcolor[rgb]{0.12,0.31,0.71}{#1}}
\newcommand{\DataTypeTok}[1]{\textcolor[rgb]{0.12,0.31,0.71}{#1}}
\newcommand{\DecValTok}[1]{\textcolor[rgb]{0.77,0.35,0.07}{#1}}
\newcommand{\DocumentationTok}[1]{\textcolor[rgb]{0.18,0.55,0.23}{#1}}
\newcommand{\ErrorTok}[1]{\textcolor[rgb]{1.00,0.00,0.00}{#1}}
\newcommand{\ExtensionTok}[1]{#1}
\newcommand{\FloatTok}[1]{\textcolor[rgb]{0.77,0.35,0.07}{#1}}
\newcommand{\FunctionTok}[1]{\textcolor[rgb]{0.44,0.19,0.63}{#1}}
\newcommand{\ImportTok}[1]{\textcolor[rgb]{0.44,0.19,0.63}{#1}}
\newcommand{\InformationTok}[1]{\textcolor[rgb]{0.38,0.63,0.69}{#1}}
\newcommand{\KeywordTok}[1]{\textcolor[rgb]{0.12,0.31,0.71}{#1}}
\newcommand{\NormalTok}[1]{#1}
\newcommand{\OperatorTok}[1]{\textcolor[rgb]{0.20,0.20,0.20}{#1}}
\newcommand{\OtherTok}[1]{\textcolor[rgb]{0.00,0.44,0.13}{#1}}
\newcommand{\PreprocessorTok}[1]{\textcolor[rgb]{0.44,0.19,0.63}{#1}}
\newcommand{\RegionMarkerTok}[1]{#1}
\newcommand{\SpecialCharTok}[1]{\textcolor[rgb]{0.63,0.19,0.43}{#1}}
\newcommand{\SpecialStringTok}[1]{\textcolor[rgb]{0.63,0.19,0.43}{#1}}
\newcommand{\StringTok}[1]{\textcolor[rgb]{0.63,0.19,0.43}{#1}}
\newcommand{\VariableTok}[1]{\textcolor[rgb]{0.10,0.09,0.49}{#1}}
\newcommand{\VerbatimStringTok}[1]{\textcolor[rgb]{0.63,0.19,0.43}{#1}}
\newcommand{\WarningTok}[1]{\textcolor[rgb]{0.38,0.63,0.69}{#1}}
\makeatletter
\@ifundefined{Shaded}{}{\renewenvironment{Shaded}{\def\FrameCommand{\fboxsep=5pt\fboxrule=0.5pt\fcolorbox{wdblue!35}{codebg}}%
   \MakeFramed{\advance\hsize-\width\FrameRestore}\small}{\endMakeFramed}}
\makeatother
\fvset{fontsize=\small}

% ---------- 表格 ----------
\renewcommand{\arraystretch}{1.42}
\newcolumntype{Y}{>{\raggedright\arraybackslash}X}

% ---------- 图片 ----------
% 课件截图做成浮动体：LaTeX 会自动用正文填满页面，避免大块空白
\setcounter{topnumber}{3}\setcounter{bottomnumber}{2}\setcounter{totalnumber}{4}
\renewcommand{\topfraction}{0.85}\renewcommand{\bottomfraction}{0.6}\renewcommand{\textfraction}{0.12}\renewcommand{\floatpagefraction}{0.75}
\setlength{\textfloatsep}{10pt plus 3pt minus 2pt}\setlength{\floatsep}{8pt plus 2pt}
\newcommand{\slidefig}[3]{\begin{figure}[!ht]\centering\begin{tcolorbox}[enhanced,width=#2\linewidth,arc=2pt,boxrule=0.6pt,colframe=wdgray!70,
   colback=white,left=2pt,right=2pt,top=2pt,bottom=2pt,drop fuzzy shadow=black!25]\includegraphics[width=\linewidth]{#1}\end{tcolorbox}
   \vspace{-3pt}{\small\color{wdgray}#3}\end{figure}}
\newcommand{\plainfig}[3]{\par\noindent\begin{minipage}{\linewidth}\centering\includegraphics[width=#2\linewidth]{#1}\par\vspace{2pt}{\small\color{wdgray}#3}\end{minipage}\par\vspace{4pt}}

\begin{document}
% ---------- 封面条（仿课件封面） ----------
\thispagestyle{plain}
\begin{tcolorbox}[enhanced,enlarge left by=-4.8cm,width=\linewidth+4.8cm,arc=0pt,boxrule=0pt,
   colback=wddark,interior style={left color=wddark,right color=wdblue},left=18pt,right=18pt,top=14pt,bottom=12pt,
   overlay={\draw[white,line width=2.5pt] ([xshift=10pt,yshift=-10pt]frame.north west) -- ++(0,-0.9cm);
            \draw[white,line width=2.5pt] ([xshift=10pt,yshift=-10pt]frame.north west) -- ++(1.6cm,0);
            \draw[white,line width=2.5pt] ([xshift=-10pt,yshift=10pt]frame.south east) -- ++(0,0.9cm);
            \draw[white,line width=2.5pt] ([xshift=-10pt,yshift=10pt]frame.south east) -- ++(-1.6cm,0);}]
  {\color{white!85!wdlightblue}\small 数据结构 · 第2章
线性表 · 康奈尔听课笔记}\par\vspace{6pt}
  {\color{white}\fontsize{22}{28}\selectfont\bfseries 2.2.1
顺序表的定义}\par\vspace{6pt}
  {\color{white!80!wdlightblue}\small 课件 2.2.1 · 正课字幕 P9（Curran
自录）· 王道书 §2.2.1（第 15--16 页）}
\end{tcolorbox}
\vspace{4pt}
\begin{tcolorbox}[enhanced,arc=3pt,boxrule=0.8pt,colframe=wdblue,colback=wdblue!3,left=10pt,right=10pt,top=6pt,bottom=6pt,
   colbacktitle=wdblue,coltitle=white,fonttitle=\bfseries\small,title={目录 · 点击标题跳转}]
\hypertarget{toc}{}\makeatletter\@starttoc{toc}\makeatother
\end{tcolorbox}
\hypertarget{ux8fd9ux4e00ux8282ux5728ux8bb2ux4ec0ux4e48}{%
\section{这一节在讲什么}\label{ux8fd9ux4e00ux8282ux5728ux8bb2ux4ec0ux4e48}}

\tocback

\Needspace{4\baselineskip}

\cuenote{2.1 讲了线性表的逻辑结构和运算，2.2 开始讲什么？}2.1
只说了线性表“逻辑上”是一串有先后的元素。从这节起回答：\textbf{在计算机里怎么存}。有两种存法，各成一种实现：

\begin{itemize}
\tightlist
\item
  顺序存储 → \textbf{顺序表}（2.2，本节先讲它的定义和代码实现）；
\item
  链式存储 → \textbf{链表}（2.3）。
\end{itemize}

下一节 2.2.2 再在顺序表上实现 2.1 定义的那些基本操作。

\slidefig{assets/2.2.1-p02.jpg}{0.6}{课件：知识总览}

\begin{bookbox}

王道书本章考点追踪：顺序表的应用是\textbf{算法题}的常客（2010、2011、2018、2020、2025
年）。这一节的结构体定义就是以后写算法题的起点，要能默写。

\end{bookbox}

\hypertarget{ux4e00ux987aux5e8fux8868ux7684ux5b9aux4e49}{%
\section{一、顺序表的定义}\label{ux4e00ux987aux5e8fux8868ux7684ux5b9aux4e49}}

\tocback

\cueblock{\cuenote}{\relax{什么是顺序表？顺序存储的关键是什么？}}

\begin{defbox}

\textbf{顺序表}：用\hlred{顺序存储}的方式实现的线性表。

\textbf{顺序存储}：把\uln{逻辑上相邻}的元素存储在\uln{物理位置上也相邻}的存储单元中，元素之间的关系由存储单元的邻接关系来体现。

\end{defbox}

王道书的说法：用一组\textbf{地址连续}的存储单元依次存储线性表中的数据元素，所以顺序表的特点是\hlred{表中元素的逻辑顺序与其物理存储顺序完全一致}。

\slidefig{assets/2.2.1-p03.jpg}{0.65}{课件：顺序表的定义}

\hypertarget{ux7b2c-i-ux4e2aux5143ux7d20ux7684ux5730ux5740ux600eux4e48ux7b97}{%
\subsection{\texorpdfstring{第 \(i\)
个元素的地址怎么算}{第 i 个元素的地址怎么算}}\label{ux7b2c-i-ux4e2aux5143ux7d20ux7684ux5730ux5740ux600eux4e48ux7b97}}

\Needspace{4\baselineskip}

\cuenote{已知首地址，第 \(i\) 个元素的地址是多少？}线性表要求元素\uln{类型相同}（2.1
的定义），所以每个元素占的空间一样大；顺序表又是连续存放，于是地址可以直接算出来。设起始地址为
\(LOC(L)\)（LOC 是 location 的缩写）：

\[LOC(a_i)=LOC(L)+(i-1)\times \texttt{sizeof(ElemType)}\]

{\rowcolors{2}{wdblue!5}{white}\arrayrulecolor{wdblue!35}
\begin{xltabular}{\linewidth}{>{\hsize=0.545\hsize\raggedright\arraybackslash}X>{\hsize=0.682\hsize\raggedright\arraybackslash}X>{\hsize=1.773\hsize\raggedright\arraybackslash}X}
\rowcolor{wdblue}\textcolor{white}{\bfseries 位序 \(i\)} & \textcolor{white}{\bfseries 数组下标} & \textcolor{white}{\bfseries 存储地址} \\
\endfirsthead
1 & 0 & \(LOC(L)\) \\ \hline
2 & 1 & \(LOC(L)+\texttt{sizeof(ElemType)}\) \\ \hline
\(i\) & \(i-1\) & \(LOC(L)+(i-1)\times\texttt{sizeof(ElemType)}\) \\ \hline
\(n\) & \(n-1\) & \(LOC(L)+(n-1)\times\texttt{sizeof(ElemType)}\) \\ \hline
\end{xltabular}\par\vspace{4pt}}

\begin{itemize}
\tightlist
\item
  \textbf{\texttt{ElemType}}：Element Type
  的缩写，就是顺序表里存的元素类型，代码里换成 \texttt{int} 或自己定义的
  \texttt{struct} 即可。
\item
  \textbf{元素大小用 \texttt{sizeof} 求}：\texttt{sizeof(int)} 通常是
  4B；下面这个 \texttt{Customer} 有两个
  \texttt{int}，\texttt{sizeof(Customer)} = 8B。不用自己算，交给
  \texttt{sizeof}。
\end{itemize}

\begin{Shaded}
\begin{Highlighting}[]
\KeywordTok{typedef} \KeywordTok{struct} \OperatorTok{\{}
    \DataTypeTok{int}\NormalTok{ num}\OperatorTok{;}      \CommentTok{// 号数}
    \DataTypeTok{int}\NormalTok{ people}\OperatorTok{;}   \CommentTok{// 人数}
\OperatorTok{\}}\NormalTok{ Customer}\OperatorTok{;}
\end{Highlighting}
\end{Shaded}

\begin{pitfallbox}

公式里是 \((i-1)\) 不是 \(i\)：位序从 1 开始，第 1
个元素就在首地址处，偏移量为 0。书上图 2.1
和注意框都在强调\uln{位序与下标之间的转换}。

\end{pitfallbox}

\hypertarget{ux4e8cux5b9eux73b0ux65b9ux5f0fux4e00ux9759ux6001ux5206ux914d}{%
\section{二、实现方式一：静态分配}\label{ux4e8cux5b9eux73b0ux65b9ux5f0fux4e00ux9759ux6001ux5206ux914d}}

\tocback

\Needspace{4\baselineskip}

\cuenote{静态分配的顺序表怎么定义？\texttt{MaxSize} 和 \texttt{length} 各管什么？}用大家最熟悉的\textbf{静态数组}存元素，数组长度一旦确定就不能变：

\begin{Shaded}
\begin{Highlighting}[]
\PreprocessorTok{\#define MaxSize }\DecValTok{10}\PreprocessorTok{              }\CommentTok{// 定义最大长度}
\KeywordTok{typedef} \KeywordTok{struct} \OperatorTok{\{}
\NormalTok{    ElemType data}\OperatorTok{[}\NormalTok{MaxSize}\OperatorTok{];}     \CommentTok{// 用静态的"数组"存放数据元素}
    \DataTypeTok{int}\NormalTok{ length}\OperatorTok{;}                 \CommentTok{// 顺序表的当前长度}
\OperatorTok{\}}\NormalTok{ SqList}\OperatorTok{;}                       \CommentTok{// 顺序表的类型定义（静态分配方式）}
\end{Highlighting}
\end{Shaded}

\begin{itemize}
\tightlist
\item
  \texttt{MaxSize}：最多能存几个元素（容量）；\texttt{length}：现在实际存了几个（表长）。
\item
  \texttt{SqList}：Sq 是 sequence（顺序）的缩写。
\end{itemize}

\hypertarget{ux521dux59cbux5316ux65f6ux5185ux5b58ux91ccux53d1ux751fux4e86ux4ec0ux4e48}{%
\subsection{初始化时内存里发生了什么}\label{ux521dux59cbux5316ux65f6ux5185ux5b58ux91ccux53d1ux751fux4e86ux4ec0ux4e48}}

老师用 \texttt{ElemType\ =\ int} 的完整代码，按执行顺序标了 ①②③：

\begin{Shaded}
\begin{Highlighting}[]
\PreprocessorTok{\#include }\ImportTok{\textless{}stdio.h\textgreater{}}
\PreprocessorTok{\#define MaxSize }\DecValTok{10}
\KeywordTok{typedef} \KeywordTok{struct} \OperatorTok{\{}
    \DataTypeTok{int}\NormalTok{ data}\OperatorTok{[}\NormalTok{MaxSize}\OperatorTok{];}
    \DataTypeTok{int}\NormalTok{ length}\OperatorTok{;}
\OperatorTok{\}}\NormalTok{ SqList}\OperatorTok{;}

\DataTypeTok{void}\NormalTok{ InitList}\OperatorTok{(}\NormalTok{SqList }\OperatorTok{\&}\NormalTok{L}\OperatorTok{)} \OperatorTok{\{}
    \ControlFlowTok{for} \OperatorTok{(}\DataTypeTok{int}\NormalTok{ i }\OperatorTok{=} \DecValTok{0}\OperatorTok{;}\NormalTok{ i }\OperatorTok{\textless{}}\NormalTok{ MaxSize}\OperatorTok{;}\NormalTok{ i}\OperatorTok{++)}
\NormalTok{        L}\OperatorTok{.}\NormalTok{data}\OperatorTok{[}\NormalTok{i}\OperatorTok{]} \OperatorTok{=} \DecValTok{0}\OperatorTok{;}          \CommentTok{// ② 所有元素设默认值 0（可省略）}
\NormalTok{    L}\OperatorTok{.}\NormalTok{length }\OperatorTok{=} \DecValTok{0}\OperatorTok{;}               \CommentTok{// ③ 初始长度为 0（不能省）}
\OperatorTok{\}}

\DataTypeTok{int}\NormalTok{ main}\OperatorTok{()} \OperatorTok{\{}
\NormalTok{    SqList L}\OperatorTok{;}                   \CommentTok{// ① 声明顺序表}
\NormalTok{    InitList}\OperatorTok{(}\NormalTok{L}\OperatorTok{);}
    \ControlFlowTok{return} \DecValTok{0}\OperatorTok{;}
\OperatorTok{\}}
\end{Highlighting}
\end{Shaded}

{\rowcolors{2}{wdblue!5}{white}\arrayrulecolor{wdblue!35}
\begin{xltabular}{\linewidth}{>{\hsize=0.578\hsize\raggedright\arraybackslash}X>{\hsize=1.422\hsize\raggedright\arraybackslash}X}
\rowcolor{wdblue}\textcolor{white}{\bfseries 步骤} & \textcolor{white}{\bfseries 内存里的变化} \\
\endfirsthead
① \texttt{SqList\ L;} & 分配 \texttt{MaxSize*sizeof(ElemType)} = 10×4 = 40B 的连续空间给 \texttt{data}，另分配 4B 给 \texttt{length} \\ \hline
② \texttt{for} 循环 & 把 \texttt{data{[}0{]}}～\texttt{data{[}9{]}} 都写成 0 \\ \hline
③ \texttt{L.length=0;} & 表长置 0，此时是空表 \\ \hline
\end{xltabular}\par\vspace{4pt}}

\slidefig{assets/2.2.1-p05.jpg}{0.65}{课件：静态分配初始化的内存示意}

\hypertarget{ux54eaux4e00ux6b65ux80fdux7701ux54eaux4e00ux6b65ux4e0dux80fdux7701}{%
\subsection{哪一步能省，哪一步不能省}\label{ux54eaux4e00ux6b65ux80fdux7701ux54eaux4e00ux6b65ux4e0dux80fdux7701}}

\Needspace{4\baselineskip}

\cuenote{初始化时，给 data 设默认值能不能省？length=0 呢？}老师把 ②
去掉，然后“违规”地把整个 \texttt{data} 数组打印出来，结果前 8 个是
0，最后两个是 \texttt{-272632568}、\texttt{32766} 这种怪数：

\slidefig{assets/2.2.1-p06.jpg}{0.65}{课件：不设默认值时的“脏数据”}

\begin{itemize}
\tightlist
\item
  \textbf{怪数从哪来}：内存里有之前遗留的\hlred{脏数据}。系统分给你这块空间，但不会帮你清空。
\item
  \textbf{② 可以省}：正规的访问只到 \texttt{length}
  为止（\texttt{i\ \textless{}\ L.length}），刚初始化时
  \texttt{length=0}，一个元素都不会去读，脏数据碰不到。按
  \texttt{MaxSize} 遍历本身就是违规的。更好的做法是通过基本操作（如
  \texttt{GetElem}）访问元素。
\item
  \textbf{③ 不能省}：\texttt{length}
  自己也躺在可能有脏数据的内存里。有同学会说“C 语言 \texttt{int} 默认是
  0”，老师说那是某些编译器帮你做的，换个编译器就不一定，所以必须手动置
  0。
\end{itemize}

王道书给的静态分配初始化，正好就只有 \texttt{L.length=0;} 这一句。

\hypertarget{ux9759ux6001ux5206ux914dux7684ux5c40ux9650}{%
\subsection{静态分配的局限}\label{ux9759ux6001ux5206ux914dux7684ux5c40ux9650}}

\cueblock{\cuenote}{\relax{静态数组存满了怎么办？一开始开很大行不行？}}

\begin{itemize}
\tightlist
\item
  \textbf{存满了}：老师原话“可以放弃治疗”。静态数组长度声明后就不能改，存储空间是静态的。
\item
  \textbf{一开始开很大}：比如开 100 万但只用 10 个，浪费内存。
\item
  书上的说法：静态分配时空间一旦占满，再插入新元素会导致\textbf{溢出}，进而可能引发程序异常。
\end{itemize}

所以要让容量可变，就得用动态分配。

\hypertarget{ux4e09ux5b9eux73b0ux65b9ux5f0fux4e8cux52a8ux6001ux5206ux914d}{%
\section{三、实现方式二：动态分配}\label{ux4e09ux5b9eux73b0ux65b9ux5f0fux4e8cux52a8ux6001ux5206ux914d}}

\tocback

\cueblock{\cuenote}{\relax{动态分配的顺序表多了哪个成员？data 变成了什么？}}

\begin{Shaded}
\begin{Highlighting}[]
\PreprocessorTok{\#define InitSize }\DecValTok{10}\PreprocessorTok{             }\CommentTok{// 顺序表的初始长度}
\KeywordTok{typedef} \KeywordTok{struct} \OperatorTok{\{}
\NormalTok{    ElemType }\OperatorTok{*}\NormalTok{data}\OperatorTok{;}             \CommentTok{// 指示动态分配数组的指针}
    \DataTypeTok{int}\NormalTok{ MaxSize}\OperatorTok{;}                \CommentTok{// 顺序表的最大容量}
    \DataTypeTok{int}\NormalTok{ length}\OperatorTok{;}                 \CommentTok{// 顺序表的当前长度}
\OperatorTok{\}}\NormalTok{ SeqList}\OperatorTok{;}                      \CommentTok{// 顺序表的类型定义（动态分配方式）}
\end{Highlighting}
\end{Shaded}

和静态版比，两处变化：\texttt{data}
从数组变成了\textbf{指针}；容量会变，所以 \texttt{MaxSize}
从宏常量变成了结构体里的\textbf{变量}。

\begin{fixbox}

\textbf{名字统一}：课件里静态版叫 \texttt{SqList}、动态版叫
\texttt{SeqList}，王道书也是这样分的；书里 \texttt{InitSize} 取
100、\texttt{MaxSize} 取 50，课件都取
10，数值无所谓。后面的笔记沿用这两个类型名。

\end{fixbox}

\hypertarget{malloc-ux548c-free}{%
\subsection{malloc 和 free}\label{malloc-ux548c-free}}

\Needspace{4\baselineskip}

\cuenote{malloc 返回什么？为什么要强制类型转换？参数是什么意思？}C
语言用 \texttt{malloc} 申请、\texttt{free}
释放一整片连续的内存空间，二者在头文件
\texttt{\textless{}stdlib.h\textgreater{}} 里：

\begin{Shaded}
\begin{Highlighting}[]
\NormalTok{L}\OperatorTok{.}\NormalTok{data }\OperatorTok{=} \OperatorTok{(}\NormalTok{ElemType }\OperatorTok{*)}\NormalTok{ malloc}\OperatorTok{(}\KeywordTok{sizeof}\OperatorTok{(}\NormalTok{ElemType}\OperatorTok{)} \OperatorTok{*}\NormalTok{ InitSize}\OperatorTok{);}
\end{Highlighting}
\end{Shaded}

{\rowcolors{2}{wdblue!5}{white}\arrayrulecolor{wdblue!35}
\begin{xltabular}{\linewidth}{>{\hsize=0.696\hsize\raggedright\arraybackslash}X>{\hsize=1.304\hsize\raggedright\arraybackslash}X}
\rowcolor{wdblue}\textcolor{white}{\bfseries 部分} & \textcolor{white}{\bfseries 含义} \\
\endfirsthead
\texttt{malloc(...)} 的参数 & 要申请多大的连续空间：一个元素的大小 × 初始个数 \\ \hline
\texttt{malloc} 的返回值 & 这片空间\textbf{起始地址}的指针 \\ \hline
\texttt{(ElemType\ *)} & 把返回的指针\hlred{强制转换}成指向你的元素类型的指针（\texttt{int} 表就写 \texttt{(int\ *)}），再赋给 \texttt{L.data} \\ \hline
\end{xltabular}\par\vspace{4pt}}

\begin{itemize}
\item
  C++ 可以用 \texttt{new}、\texttt{delete}
  做同样的事，课程照顾跨考同学，主要用
  \texttt{malloc}/\texttt{free}。王道书给的 C++ 写法是：

\begin{Shaded}
\begin{Highlighting}[]
\NormalTok{L}\OperatorTok{.}\NormalTok{data }\OperatorTok{=}\NormalTok{ new ElemType}\OperatorTok{[}\NormalTok{InitSize}\OperatorTok{];}
\end{Highlighting}
\end{Shaded}
\item
  强制转换为什么不能写错，2.2.2
  顺序表的查找一节会讲透：按下标取元素时一次取几个字节，就是由指针类型决定的。
\end{itemize}

\hypertarget{ux6269ux5bb9-increasesizeux4e00ux6b65ux4e00ux6b65ux770b}{%
\subsection{扩容
IncreaseSize：一步一步看}\label{ux6269ux5bb9-increasesizeux4e00ux6b65ux4e00ux6b65ux770b}}

\cueblock{\cuenote}{\relax{动态顺序表满了怎么扩容？为什么要先存一个指针 p？}}

\begin{Shaded}
\begin{Highlighting}[]
\PreprocessorTok{\#include }\ImportTok{\textless{}stdlib.h\textgreater{}}
\DataTypeTok{void}\NormalTok{ InitList}\OperatorTok{(}\NormalTok{SeqList }\OperatorTok{\&}\NormalTok{L}\OperatorTok{)} \OperatorTok{\{}
\NormalTok{    L}\OperatorTok{.}\NormalTok{data }\OperatorTok{=} \OperatorTok{(}\DataTypeTok{int} \OperatorTok{*)}\NormalTok{ malloc}\OperatorTok{(}\NormalTok{InitSize }\OperatorTok{*} \KeywordTok{sizeof}\OperatorTok{(}\DataTypeTok{int}\OperatorTok{));}  \CommentTok{// 申请 10 个 int 的连续空间}
\NormalTok{    L}\OperatorTok{.}\NormalTok{length }\OperatorTok{=} \DecValTok{0}\OperatorTok{;}
\NormalTok{    L}\OperatorTok{.}\NormalTok{MaxSize }\OperatorTok{=}\NormalTok{ InitSize}\OperatorTok{;}
\OperatorTok{\}}

\DataTypeTok{void}\NormalTok{ IncreaseSize}\OperatorTok{(}\NormalTok{SeqList }\OperatorTok{\&}\NormalTok{L}\OperatorTok{,} \DataTypeTok{int}\NormalTok{ len}\OperatorTok{)} \OperatorTok{\{}  \CommentTok{// 增加动态数组的长度}
    \DataTypeTok{int} \OperatorTok{*}\NormalTok{p }\OperatorTok{=}\NormalTok{ L}\OperatorTok{.}\NormalTok{data}\OperatorTok{;}                       \CommentTok{// ① p 记住旧空间}
\NormalTok{    L}\OperatorTok{.}\NormalTok{data }\OperatorTok{=} \OperatorTok{(}\DataTypeTok{int} \OperatorTok{*)}\NormalTok{ malloc}\OperatorTok{((}\NormalTok{L}\OperatorTok{.}\NormalTok{MaxSize }\OperatorTok{+}\NormalTok{ len}\OperatorTok{)} \OperatorTok{*} \KeywordTok{sizeof}\OperatorTok{(}\DataTypeTok{int}\OperatorTok{));}  \CommentTok{// ② 申请更大的新空间}
    \ControlFlowTok{for} \OperatorTok{(}\DataTypeTok{int}\NormalTok{ i }\OperatorTok{=} \DecValTok{0}\OperatorTok{;}\NormalTok{ i }\OperatorTok{\textless{}}\NormalTok{ L}\OperatorTok{.}\NormalTok{length}\OperatorTok{;}\NormalTok{ i}\OperatorTok{++)}
\NormalTok{        L}\OperatorTok{.}\NormalTok{data}\OperatorTok{[}\NormalTok{i}\OperatorTok{]} \OperatorTok{=}\NormalTok{ p}\OperatorTok{[}\NormalTok{i}\OperatorTok{];}                  \CommentTok{// ③ 把旧数据复制到新区域}
\NormalTok{    L}\OperatorTok{.}\NormalTok{MaxSize }\OperatorTok{=}\NormalTok{ L}\OperatorTok{.}\NormalTok{MaxSize }\OperatorTok{+}\NormalTok{ len}\OperatorTok{;}           \CommentTok{// ④ 最大容量增加 len}
\NormalTok{    free}\OperatorTok{(}\NormalTok{p}\OperatorTok{);}                               \CommentTok{// ⑤ 释放旧空间}
\OperatorTok{\}}

\DataTypeTok{int}\NormalTok{ main}\OperatorTok{()} \OperatorTok{\{}
\NormalTok{    SeqList L}\OperatorTok{;}
\NormalTok{    InitList}\OperatorTok{(}\NormalTok{L}\OperatorTok{);}
    \CommentTok{// ...往顺序表中插满 10 个元素...}
\NormalTok{    IncreaseSize}\OperatorTok{(}\NormalTok{L}\OperatorTok{,} \DecValTok{5}\OperatorTok{);}
    \ControlFlowTok{return} \DecValTok{0}\OperatorTok{;}
\OperatorTok{\}}
\end{Highlighting}
\end{Shaded}

以表已存满（\texttt{length=10}、\texttt{MaxSize=10}）后调用
\texttt{IncreaseSize(L,\ 5)} 为例：

{\rowcolors{2}{wdblue!5}{white}\arrayrulecolor{wdblue!35}
\begin{xltabular}{\linewidth}{>{\hsize=0.500\hsize\raggedright\arraybackslash}X>{\hsize=1.600\hsize\raggedright\arraybackslash}X>{\hsize=1.400\hsize\raggedright\arraybackslash}X>{\hsize=0.800\hsize\raggedright\arraybackslash}X>{\hsize=0.700\hsize\raggedright\arraybackslash}X}
\rowcolor{wdblue}\textcolor{white}{\bfseries 步} & \textcolor{white}{\bfseries 执行} & \textcolor{white}{\bfseries \texttt{L.data} 指向} & \textcolor{white}{\bfseries \texttt{p} 指向} & \textcolor{white}{\bfseries \texttt{MaxSize}} \\
\endfirsthead
前 & 表已满 & 旧空间（10 格，全满） & --- & 10 \\ \hline
① & \texttt{p\ =\ L.data} & 旧空间 & 旧空间 & 10 \\ \hline
② & \texttt{malloc} 15 格 & \textbf{新空间}（15 格，还是空的） & 旧空间 & 10 \\ \hline
③ & \texttt{for} 复制 10 次 & 新空间，前 10 格有数据 & 旧空间 & 10 \\ \hline
④ & \texttt{MaxSize\ +=\ 5} & 新空间 & 旧空间 & 15 \\ \hline
⑤ & \texttt{free(p)} & 新空间 & 旧空间已还给系统 & 15 \\ \hline
\end{xltabular}\par\vspace{4pt}}

\slidefig{assets/2.2.1-p09.jpg}{0.7}{课件：动态分配与扩容的内存过程}

\begin{itemize}
\tightlist
\item
  \textbf{为什么要先用 \texttt{p} 记住旧地址}：② 执行后 \texttt{L.data}
  就指向新空间了，如果没有 \texttt{p}，旧空间的地址就丢了，第 ③
  步没法复制，第 ⑤ 步也没法释放（这片内存就再也回收不了）。
\item
  \textbf{\texttt{p}
  本身}：是函数里的局部变量，函数结束时系统自动回收它占的那几个字节；\texttt{free(p)}
  释放的是 \texttt{p} 指向的那 10 格旧空间。
\item
  \textbf{代价}：③ 要把 \(n\) 个元素逐个复制，\hlred{时间开销大}，是
  \(O(n)\)。
\item
  \texttt{realloc}
  也能做这件事，但老师说它调用时有意想不到的坑，建议初学者用
  \texttt{malloc} + \texttt{free}，也更能看清背后的过程。
\end{itemize}

\begin{pitfallbox}

王道书注意：\hlred{动态分配不是链式存储}，它同样属于顺序存储结构，物理结构没有变化，依然支持随机存取，只是存储空间的大小可以在运行时动态调整。

\end{pitfallbox}

\hypertarget{ux56dbux987aux5e8fux8868ux7684ux7279ux70b9}{%
\section{四、顺序表的特点}\label{ux56dbux987aux5e8fux8868ux7684ux7279ux70b9}}

\tocback

\Needspace{4\baselineskip}

\cuenote{顺序表有哪些优缺点？各自的原因是什么？}不管静态还是动态分配，顺序表都有下面这些特点（课件列四条，王道书归成两优两缺）：

{\rowcolors{2}{wdblue!5}{white}\arrayrulecolor{wdblue!35}
\begin{xltabular}{\linewidth}{>{\hsize=0.830\hsize\raggedright\arraybackslash}X>{\hsize=1.787\hsize\raggedright\arraybackslash}X>{\hsize=0.383\hsize\raggedright\arraybackslash}X}
\rowcolor{wdblue}\textcolor{white}{\bfseries 特点} & \textcolor{white}{\bfseries 原因} & \textcolor{white}{\bfseries 书上} \\
\endfirsthead
① \hlred{随机存取}：\(O(1)\) 时间找到第 \(i\) 个元素，代码就是 \texttt{data{[}i-1{]}} & 元素连续存放、大小相同，首地址加偏移量直接算出地址 & 优点 \\ \hline
② \textbf{存储密度高}：每个结点只存数据元素本身 & 链式存储还要额外存指针 & 优点 \\ \hline
③ 拓展容量不方便 & 静态分配不能扩；动态分配要整体复制，时间复杂度高 & 缺点（书：要求分配连续的存储空间，不够灵活） \\ \hline
④ 插入、删除不方便，要移动大量元素 & 要保持“物理相邻”，中间插一个、删一个，后面都得挪 & 缺点 \\ \hline
\end{xltabular}\par\vspace{4pt}}

\begin{fixbox}

\textbf{名字统一}：课件和老师说“随机访问”，王道书两种都用（定义处说“随机存取”，优点处说“随机访问”）。考题里常见的是\textbf{随机存取}，本笔记统一用它，两者是一个意思。

\end{fixbox}

\begin{itemize}
\tightlist
\item
  \textbf{存储密度} = 结点中数据元素本身占的空间 ÷
  结点占的总空间。顺序表的结点只存数据，存储密度为
  1；链表的结点还要存指针，小于 1。
\item
  第 ④ 条下一节会用代码和时间复杂度具体体会。
\end{itemize}

\slidefig{assets/2.2.1-p10.jpg}{0.65}{课件：顺序表的特点}

\hypertarget{ux4e94ux77e5ux8bc6ux56deux987e}{%
\section{五、知识回顾}\label{ux4e94ux77e5ux8bc6ux56deux987e}}

\tocback

\slidefig{assets/2.2.1-p11.jpg}{0.65}{课件：知识回顾与重要考点}

\begin{summarybox}

\textbf{顺序表 =
用顺序存储实现的线性表}：逻辑上相邻的元素物理上也相邻，所以第 \(i\)
个元素的地址是
\(LOC(L)+(i-1)\times\texttt{sizeof(ElemType)}\)，\hlred{随机存取}，\(O(1)\)。

实现有两种。\textbf{静态分配}用定长数组 \texttt{data{[}MaxSize{]}} 加
\texttt{length}，初始化时 \texttt{length=0}
必须做（防脏数据），给元素设默认值可以省；满了就没办法。\textbf{动态分配}用指针
\texttt{data} 加 \texttt{MaxSize}、\texttt{length}，\texttt{malloc}
申请、\texttt{free} 释放，返回的指针要强转成元素类型指针；扩容时先用
\texttt{p} 记住旧空间，申请更大的新空间、复制、改
\texttt{MaxSize}、释放旧空间，复制带来 \(O(n)\)
的开销。动态分配仍是顺序存储，不是链式存储。

优点：随机存取、存储密度高；缺点：扩容不方便、插入删除要移动大量元素。

\textbf{自测}：① 每个元素占 4B，首地址 2000，第 5 个元素的地址是多少？②
\texttt{InitList} 里哪一步不能省？为什么？③ \texttt{IncreaseSize}
里如果不先写 \texttt{int\ *p\ =\ L.data;} 会出什么问题？④
动态分配的顺序表是链式存储吗？

\end{summarybox}
\end{document}
